\documentclass[11pt]{article}

\usepackage[utf8]{inputenc}
\usepackage[T1]{fontenc}
\usepackage[letterpaper,margin=1in]{geometry}
\usepackage{enumitem}
\usepackage{parskip}
\usepackage{hyperref}
\hypersetup{colorlinks=true,linkcolor=blue,urlcolor=blue}

\title{\textbf{Senior Project Proposal}\\[0.35em]
\large SimplrAI: An AI Explanation System for Student Learning}
\author{Soumil Bhandari}
\date{March 29, 2026}

\begin{document}
\maketitle

\section*{Introduction}

Over the last two years, artificial intelligence has moved rapidly from a
futuristic concept to an everyday academic tool. Students now routinely turn
to AI models for help with homework, reading comprehension, math explanations,
and even research summaries. According to the Pew Research Center (2024), over
63\% of American teens report using AI for schoolwork, and nearly 20\% use it
daily. Yet despite AI's explosive adoption, most tools used by students---such
as ChatGPT, Google Gemini, and Perplexity---were not intentionally designed for
pedagogical clarity, grade-level appropriateness, or cognitive simplicity.
Their outputs can be verbose, overly complex, disorganized, or subtly
inaccurate. Paradoxically, a tool intended to simplify learning can
unintentionally increase confusion.

This project responds directly to that gap. My research question is:

\begin{quote}
\itshape
How can an AI-based academic explanation system be engineered to produce fast,
accurate, and grade-appropriate explanations that reduce cognitive load and
meaningfully improve student comprehension?
\end{quote}

From this question emerges an applied research problem: identifying the optimal
combination of model architecture, prompt-engineering strategy, and
user-experience design that yields the clearest and most educationally
effective explanations for students from middle school through college.

This proposal outlines the theoretical foundations, prior research, experimental
plan, analysis approach, and projected outcomes necessary to evaluate and build
SimplrAI, an educational explanation tool designed around clarity, structure,
and learning outcomes, which transcends the raw generation capabilities of
current AI tools.

\section*{Background}

My interest in this project emerges from years of direct experience both as a
student navigating advanced coursework and as a peer tutor helping others
interpret difficult material. Throughout classes like AP Calculus BC, AP
Computer Science, Multivariable Calculus, and Linear Algebra, I often used large
language models to clarify difficult concepts. What I found was consistent: the
answers were longer than necessary, poorly structured, and occasionally subtly
wrong. Even when accurate, the explanations were not optimized for how students
actually learn.

As a debate coach and academic tutor, I witnessed a parallel issue. Students
frequently used AI tools to ``understand'' complicated topics, but the responses
they received overwhelmed them with unnecessary detail. They often ended up more
confused than before. For middle school students, the problem was even more
pronounced: explanations were too advanced, written at a collegiate reading
level, or lacked intuitive examples and analogies.

These experiences sparked my earlier project, Notes360.ai, which focused on
simplifying and structuring class notes. Working on that system taught me
important lessons about user-centered design, cognitive load, and how students
interact with educational technology. It also gave me practical experience in
full-stack engineering: React, Supabase, backend APIs, and LLM integration.

Meanwhile, my academic interest in AI deepened through my research internships,
robotics programming, and coursework in Data Structures \& Algorithms,
Differential Equations, and computer science theory. I built multiple small
LLM-based tools, experimented with prompt templates, and studied model behavior
across different architectures.

These experiences reinforced my belief that the future of educational AI depends
not just on generative power, but on design philosophy---whether the tool is
built to mirror human teaching, simplify complex ideas, and adapt explanations
to the learner's level.

SimplrAI represents the culmination of these interests: an applied research
project that combines computer science, AI engineering, cognitive science,
educational psychology, and user-experience design. It allows me to both build
a system and study its impact.

\section*{Literature Review}

\subsection*{1. AI in Education: Promise and Problems}

Since 2020, AI-based tutoring systems have gained interest as potential
supplements to human instruction. Koedinger et al. (2022) found that AI tutors
can support conceptual understanding when explanations are structured and
scaffolded. However, most commercially available LLMs lack intentional
pedagogical design. OpenAI in 2023 and Anthropic in 2024 both note that their
models are optimized for general reasoning, not grade-level learning.

Pew Research in 2024 confirms rapid adoption among students, but warns that
``AI's educational usefulness is limited by inconsistent accuracy and lack of
transparent reasoning.'' These concerns align with findings from Dwivedi et al.,
who argue that LLMs produce ``syntactically fluent but epistemically
unreliable'' outputs when used for learning.

Thus, the field acknowledges the strong potential of AI tutors but highlights an
unresolved problem: LLMs are not built for education, and students require
structured guidance.

\subsection*{2. Cognitive Load Theory: Why Structure Matters}

Cognitive load theory (Sweller, 1988) establishes that learners can only process
a limited amount of novel information at once. Explanations that are too long,
too complex, or poorly organized increase extraneous load, reducing
comprehension. Later expansions of the theory (Kalyuga, 2009; Paas \& van
Merri\"enboer, 2020) show that explanations must be concise, structured, and
sequenced appropriately to maximize retention.

Chi and Wylie's ICAP framework (2014) categorizes learning activities by
engagement level, finding that structured explanations significantly improve
comprehension compared to passive reception.

This literature suggests that AI explanations require intentional
formatting---definitions, examples, analogies, and applications---to support
comprehension.

\subsection*{3. Hallucinations and Reliability in Large Language Models}

A major barrier to AI in education is factual unreliability. Maynez et al. (2020)
show that LLMs often hallucinate facts, especially when prompted broadly.
OpenAI's technical report (2023) confirms that prompting strategies dramatically
influence accuracy.

Several studies demonstrate that structured prompt templates (Li et al., 2023),
retrieval-augmented generation (Lewis et al., 2020), and constrained decoding
reduce hallucinations. Yet few studies evaluate these methods specifically for
pedagogical content generation, leaving an unexplored research gap.

\subsection*{4. UX and Learning Interfaces}

Educational technology research emphasizes the importance of interface clarity.
Nielsen Norman Group (2023) found that students learn best from interfaces with
low visual noise, predictable patterns, and progressive disclosure.

Studies of EdTech apps such as Khan Academy and Duolingo show that simple,
chunked text improves retention (Reich, 2020). Conversely, Perplexity and
ChatGPT, though technically powerful, deliver long, dense paragraphs, which
Quinn et al. (2023) identify as suboptimal for K--12 learners.

The literature reveals a UX gap: AI tools are not designed around the cognitive
needs of younger learners.

\subsection*{5. Grade-Level Appropriateness and Reading Complexity}

Text complexity research demonstrates the importance of tailoring explanations
to student reading levels. The Lexile Framework (MetaMetrics, 2019) shows that
comprehension drops sharply when text exceeds a learner's threshold. Shanahan
(2014) argues that academic content must be adapted---not merely
simplified---to align with cognitive development stages.

AI models rarely do this automatically. Studies by Xu et al. (2023) show that
LLMs default to college-level prose unless explicitly constrained. This
reinforces the need for grade-specific prompting, one of SimplrAI's core
features.

\subsection*{6. The Gap in the Literature}

Across AI education research, three consistent gaps emerge:

\begin{itemize}[leftmargin=1.6em]
  \item Lack of grade-appropriate explanation systems
  \item Lack of research testing structured prompt templates for student
  comprehension
  \item Lack of user-experience studies evaluating how students interact with
  AI explanations
\end{itemize}

These gaps define the contribution of SimplrAI.

\section*{Significance}

Given AI's increasing role in education, developing systems that actually
enhance student comprehension is critical. Current LLMs provide powerful but
unstructured responses, often mismatched to student reading levels. By creating
and evaluating SimplrAI, this research fills a pressing need: an AI system that
prioritizes clarity, accuracy, and cognitive efficiency.

My findings will contribute to:

\begin{itemize}[leftmargin=1.6em]
  \item Human--AI interaction research, showing how users interpret AI
  explanations
  \item Educational psychology, providing new data on structured explanations
  \item AI safety research, evaluating hallucination-reduction methods in
  academic contexts
  \item UX research, identifying interface patterns that improve learning
  outcomes
\end{itemize}

Additionally, my results may directly benefit teachers and tutoring programs
seeking tools that help rather than overwhelm students.

Most importantly, this research addresses a real social need: ensuring that the
AI tools students increasingly rely on actually help them learn.

\section*{Method}

\subsection*{A. Experimental Phase}

\textbf{1. Information Needed}

To answer my research question, I need:

\begin{itemize}[leftmargin=1.6em]
  \item Explanations generated by SimplrAI
  \item User feedback from students evaluating clarity and comprehension
  \item Quantitative scores based on structured rubrics
  \item Reading-level metrics for each explanation
  \item Accuracy assessments by subject-area experts (teachers or advanced
  students)
\end{itemize}

\textbf{2. Participants}

I will recruit middle school, high school, and college students (ages 11--22).
All participants will receive:

\begin{itemize}[leftmargin=1.6em]
  \item Full informed consent (and parental consent for minors)
  \item Anonymity
  \item The right to withdraw at any time
\end{itemize}

No identifying data will be collected.

\textbf{3. Materials}

\begin{itemize}[leftmargin=1.6em]
  \item SimplrAI prototype
  \item Surveys delivered via Google Forms
  \item Explanation comparison worksheets
  \item Reading-level measurement tools (Lexile, Coh-Metrix, GPT-based
  complexity scoring)
\end{itemize}

\textbf{4. Procedure}

Participants in various subjects and grade levels will use SimplrAI before their
quiz or tests. They will take the test, and complete a questionnaire which will
collect personal information like:

\begin{itemize}[leftmargin=1.6em]
  \item Age
  \item Gender
  \item Familiarity with AI
\end{itemize}

They rate both on:

\begin{itemize}[leftmargin=1.6em]
  \item Accuracy
  \item Clarity
  \item Helpfulness
  \item Reading difficulty
  \item Trust
\end{itemize}

They then report their score on the examination. Data is collected and
anonymized.

\subsection*{B. Ethical Protections}

\begin{itemize}[leftmargin=1.6em]
  \item Minors will require parental consent
  \item All surveys will omit identifying information
  \item No sensitive topics will be included
  \item Participants will not be paid, ensuring no coercion
  \item All data stored securely in encrypted files
\end{itemize}

\subsection*{C. Analysis Phase}

I will use:

\begin{itemize}[leftmargin=1.6em]
  \item Descriptive statistics (mean, SD, distributions)
  \item Paired t-tests comparing SimplrAI vs ChatGPT clarity scores
  \item Content analysis of open-ended responses
  \item Accuracy scoring using expert-review rubrics
\end{itemize}

\subsection*{D. Challenges}

\begin{itemize}[leftmargin=1.6em]
  \item Recruiting enough participants --- I will partner with tutoring centers
  and online student groups
  \item Ensuring test neutrality --- prompts will be controlled
  \item Model API rate limits --- I will cache responses
  \item Variability in student background --- sample will be stratified
\end{itemize}

\section*{Final Product and Value}

The final product of this project will be a comprehensive, research-driven suite
of deliverables centered around SimplrAI, designed to bridge theoretical
research on AI explanations with practical classroom use. At its core will be a
fully functional SimplrAI web application, allowing users to interact with
AI-generated explanations, compare different prompting and explanation
strategies, and evaluate their effectiveness across multiple academic domains.
The platform will be designed with accessibility and usability in mind, making
it suitable for both students and educators without requiring advanced technical
background.

Accompanying the application will be a 20--25 page formal research paper that
rigorously documents the project's motivation, methodology, experimental design,
evaluation metrics, results, and limitations. This paper will situate the work
within existing literature on explainable AI, learning sciences, and human-AI
interaction, while clearly articulating the novel contributions of this project.
Detailed appendices will provide transparency and reproducibility, including
curated data tables, system screenshots, prompt templates, and evaluation
rubrics used throughout the study.

In addition, the project will feature a public explainer website that translates
the technical findings into a clear, engaging narrative for a broader audience.
This site will summarize key insights, demonstrate the tool's real-world
applications, and showcase examples of how SimplrAI can support learning,
critical thinking, and conceptual understanding in educational settings.

Collectively, this product contributes a replicable and well-documented
methodology for evaluating AI explanations, addressing a gap between abstract
explainability research and classroom-ready tools. Beyond its academic value,
it delivers a practical, research-validated resource that students and teachers
can immediately use to better understand not just what an AI says, but why and
how those explanations support learning.

\section*{Timeline and Work Schedule}

\subsection*{January --- Pre-Internship: Planning \& Setup}

\textit{Jan 5 -- Jan 15}
\begin{itemize}[leftmargin=1.6em,nosep]
  \item Finalize research question and hypothesis
  \item Lock evaluation metrics (clarity, accuracy, reading level, trust)
  \item Draft informed consent forms and surveys
  \item Identify comparison baseline (ChatGPT default explanations)
\end{itemize}

\textit{Jan 16 -- Jan 23}
\begin{itemize}[leftmargin=1.6em,nosep]
  \item Submit full research proposal for approval
  \item Begin SimplrAI system architecture planning (UX layout + prompt
  framework)
  \item Set up data collection pipeline (Google Forms, spreadsheets)
\end{itemize}

\subsection*{February --- Prototype Development \& Pilot Testing}

\textit{Feb 1 -- Feb 15}
\begin{itemize}[leftmargin=1.6em,nosep]
  \item Build SimplrAI MVP (frontend + LLM backend)
  \item Implement grade-level prompting and structured explanation templates
  \item Integrate reading-level scoring tools (Lexile / Coh-Metrix / GPT-based)
\end{itemize}

\textit{Feb 16 -- Feb 28}
\begin{itemize}[leftmargin=1.6em,nosep]
  \item Pilot test with small sample (5--10 students)
  \item Refine prompts and UI based on feedback
  \item Finalize experimental procedure
\end{itemize}

\subsection*{March --- Onsite Placement Begins + Formal Data Collection}

\textit{Mar 2 -- Mar 6}
\begin{itemize}[leftmargin=1.6em,nosep]
  \item Start onsite placement
\end{itemize}

\textit{Mar 7 -- Mar 20}
\begin{itemize}[leftmargin=1.6em,nosep]
  \item Full-scale participant recruitment
  \item Collect explanation ratings and test score data
  \item Conduct accuracy reviews with advanced students/teachers
\end{itemize}

\textit{Mar 27} --- Paper sections due: Introduction, Literature Review, Method.

\subsection*{April --- Analysis \& Writing Phase}

\textit{Apr 1 -- Apr 15}
\begin{itemize}[leftmargin=1.6em,nosep]
  \item Clean and anonymize dataset
  \item Run statistical analysis (means, SDs, paired t-tests)
  \item Analyze open-ended responses
\end{itemize}

\textit{Apr 16 -- Apr 30}
\begin{itemize}[leftmargin=1.6em,nosep]
  \item Write Results and Discussion sections
  \item Capture screenshots and system diagrams
  \item Prepare appendix materials (surveys, rubrics, consent forms)
\end{itemize}

\subsection*{May --- Final Product \& Presentation}

\textit{May 1} --- Draft 1 of Final Product Due.

\textit{May 2 -- May 14}
\begin{itemize}[leftmargin=1.6em,nosep]
  \item Revise paper based on feedback
  \item Finalize SimplrAI demo and explainer website
\end{itemize}

\textit{May 15} --- Final Draft Due / End of Onsite Placement.

\textit{May 28 -- May 29} --- Senior Showcase presentation.

\section*{Bibliography}

\textbf{AI / LLM Research}

\begin{list}{}{%
  \setlength{\leftmargin}{1.5em}
  \setlength{\itemindent}{-1.5em}
  \setlength{\itemsep}{0.3em}
}
\item OpenAI. (2023). \textit{GPT-4 Technical Report.}
\url{https://arxiv.org/abs/2303.08774}
\item Anthropic. (2024). \textit{Constitutional AI: Harmlessness from AI
Feedback.} \url{https://www.anthropic.com/index/constitutional-ai}
\item Li, X., Liu, H., \& Tsvetkov, Y. (2023). Prompt Engineering for Large
Language Models: A Survey. \url{https://arxiv.org/abs/2303.08774}
\item Maynez, J., Narayan, S., Bohnet, B., \& McDonald, R. (2020). On
Faithfulness and Factuality in Abstractive Summarization.
\url{https://aclanthology.org/2020.acl-main.450/}
\item Xu, W., Li, Z., \& Zhao, T. (2023). Grade-level Text Controls for LLMs.
\url{https://arxiv.org/abs/2310.10572}
\end{list}

\textbf{Cognitive Load / Educational Psychology}

\begin{list}{}{%
  \setlength{\leftmargin}{1.5em}
  \setlength{\itemindent}{-1.5em}
  \setlength{\itemsep}{0.3em}
}
\item Sweller, J. (1988). Cognitive load during problem solving: Effects on
learning. \textit{Cognitive Science}, 12(2), 257--285.
\url{https://doi.org/10.1207/s15516709cog1202_4}
\item Kalyuga, S. (2009). \textit{Managing Cognitive Load in Adaptive Learning
Environments.} \url{https://link.springer.com/book/10.1007/978-1-4419-8126-4}
\item Paas, F., \& van Merri\"enboer, J. (2020). Cognitive-load theory: A review
of recent developments. \url{https://doi.org/10.1007/s10648-020-09534-9}
\item Chi, M. T. H., \& Wylie, R. (2014). The ICAP Framework: Linking cognitive
engagement to active learning outcomes.
\url{https://doi.org/10.3102/0034654314542828}
\end{list}

\textbf{EdTech, UX, Interface Design}

\begin{list}{}{%
  \setlength{\leftmargin}{1.5em}
  \setlength{\itemindent}{-1.5em}
  \setlength{\itemsep}{0.3em}
}
\item Nielsen Norman Group. (2023). \textit{10 Usability Heuristics for User
Interface Design.} \url{https://www.nngroup.com/articles/ten-usability-heuristics/}
\item Reich, J. (2020). \textit{Failure to Disrupt: Why Technology Alone Can't
Transform Education.} MIT Press.
\url{https://mitpress.mit.edu/9780262539735/failure-to-disrupt/}
\item Quinn, D., Cooc, N., \& McIntyre, J. (2023). AI Tools and K--12 Learning:
A Review of Opportunities and Risks.
\url{https://journals.sagepub.com/doi/full/10.3102/00346543231160869}
\end{list}

\textbf{Reading Difficulty / Grade Level Research}

\begin{list}{}{%
  \setlength{\leftmargin}{1.5em}
  \setlength{\itemindent}{-1.5em}
  \setlength{\itemsep}{0.3em}
}
\item MetaMetrics. (2019). \textit{Lexile Framework for Reading.}
\url{https://lexile.com/}
\item Shanahan, T. (2014). The complexity of text and the common core.
\url{https://www.shanahanonliteracy.com/blog/the-complexity-of-text-and-the-common-core}
\end{list}

\textbf{AI in Education, Adoption, Social Trends}

\begin{list}{}{%
  \setlength{\leftmargin}{1.5em}
  \setlength{\itemindent}{-1.5em}
  \setlength{\itemsep}{0.3em}
}
\item Pew Research Center. (2024). \textit{Teens and AI: How Youth Are Using New
Technologies.} \url{https://www.pewresearch.org/internet/2024/05/01/teens-and-ai/}
\item Koedinger, K. R., et al. (2022). Learning science and AI: Using learning
engineering to design AI-powered education.
\url{https://www.pnas.org/doi/10.1073/pnas.2122730119}
\item Dwivedi, Y., et al. (2023). Theoretical Lens to Understand the Risks and
Ethics of Generative AI. \url{https://doi.org/10.1016/j.ijinfomgt.2023.102695}
\item Lewis, P., Perez, E., et al. (2020). Retrieval-Augmented Generation for
Knowledge-Intensive NLP. \url{https://arxiv.org/abs/2005.11401}
\end{list}

\end{document}
